\chapter{競合リスク}\label{ch:11}

\begin{goalbox}
\begin{itemize}
\item 競合リスク（competing risks）の状況を識別し，競合イベントと打ち切りの違いを説明できる．
\item 原因別ハザード $\lambda_j(t)$，部分密度 $f_j(t)$，累積発生関数（CIF）$F_j(t)$ を定義し，$\lambda=\sum_j\lambda_j$，$f_j=\lambda_jS$，$F_j(t)=\int_0^t\lambda_jS$ を導出できる．
\item $F_j(\infty)=\Prob(J=j)<1$ であること，$1-F_j(t)>S_j(t)$ を示せる．
\item Aalen--Johansen 型の推定量で CIF を手計算できる．
\item 競合イベントを打ち切りとみなした $1-\mathrm{KM}$ が CIF を過大推定することを示し，理由を説明できる．
\item 原因別ハザードの Cox モデルと Fine--Gray モデル（部分分布ハザード）の解釈の違いを説明できる．
\end{itemize}
\end{goalbox}

\begin{exambox}
\begin{itemize}
\item \kakomon{2021}{1}〔3〕〜〔5〕：心血管死（イベント1）と他因死（イベント2）．$F_j(t)=\int_0^t\lambda_j(u)S_2(u)\dd u$ の導出（試験の $S_2$ は本章の全生存関数 $S$），
      $\hat F_j(t)=\sum_{t_{ji}\le t}\frac{d_{ji}}{n_{ji}}\hat S_2(t_{ji}^-)$ による推定表（$\hat F_1(18)=0.75$，$\hat F_2=0.25$），
      $1-\hat S_1(t)$（他因死を打ち切り扱い）との比較と過大推定の理由．
\item \kakomon{2025}{3}：$m$ 種類の競合する死因．$\lambda_j=f_j/S$，$1-F_j(t)>S_j(t)$，$\lambda=\sum\lambda_j$ の証明，
      Weibull 型原因別ハザードの下での $S(t)=\exp\{-(\gamma_1t)^{\alpha_1}-(\gamma_2t)^{\alpha_2}\}$，
      原因別ハザードの比例ハザードモデルの $\beta$ で $F_1(t)$ への効果を評価することの適否．
\end{itemize}
\end{exambox}

\flowline{時間＋\\イベントの種類}{CIF $F_j(t)$\\原因別ハザード}{Aalen--Johansen\\Fine--Gray}{独立打ち切り\\（競合は打ち切りでない）}{$\sum\frac{d_j}{n}\hat S(t^-)$}{$1-$KM の過大推定}

\section{問題設定}\label{sec:11-setting}
がん死を主要評価項目とする研究で，他因死が起こると以後がん死は観察できない．
このように，あるイベントの発生によって別のイベントが起こりえなくなる（あるいは意味が変わる）とき，それらを\emphx{競合リスク}という．
観測されるのは最初のイベントまでの時間 $T$ と，その種類 $J\in\{1,\dots,m\}$，および打ち切り時間 $C$ である．
\[
Y=\min(T,C),\qquad J^*=\begin{cases}J&(T\le C)\\0&(T>C)\ \text{（打ち切り）}\end{cases}
\]
\Youchui 打ち切りは「イベントはいずれ起こるが観察できなかった」，競合イベントは「そのイベントは（少なくとも最初のイベントとしては）二度と起こらない」である．
両者を同じに扱うことが本章で扱う誤りの根源である．

\section{基本概念}\label{sec:11-basic}
\Hissu
\begin{teigi}[競合リスクの関数]\label{def:11-functions}
\begin{align*}
&\text{全生存関数} & S(t)&=\Prob(T>t),\\
&\text{原因別ハザード} & \lambda_j(t)&=\lim_{\Delta\to0}\frac{\Prob(t\le T<t+\Delta,\ J=j\mid T\ge t)}{\Delta},\\
&\text{部分密度} & f_j(t)&=\lim_{\Delta\to0}\frac{\Prob(t\le T<t+\Delta,\ J=j)}{\Delta},\\
&\text{累積発生関数（CIF）} & F_j(t)&=\Prob(T\le t,\ J=j)=\int_0^tf_j(u)\,\dd u.
\end{align*}
\end{teigi}
原因別ハザードは「時点 $t$ までにどのイベントも起こしていない人が，次の瞬間に原因 $j$ のイベントを起こす率」である．

\begin{teiri}[基本関係式]\label{thm:11-relations}
\begin{align}
\lambda_j(t)&=\frac{f_j(t)}{S(t)},\qquad f_j(t)=\lambda_j(t)S(t),\label{eq:11-fj}\\
\lambda(t)&=\sum_{j=1}^m\lambda_j(t),\qquad S(t)=\exp\Bigl\{-\sum_{j=1}^m\Lambda_j(t)\Bigr\},\ \ \Lambda_j(t)=\int_0^t\lambda_j,\label{eq:11-sum}\\
F_j(t)&=\int_0^t\lambda_j(u)S(u)\,\dd u,\qquad \sum_{j=1}^mF_j(t)=1-S(t).\label{eq:11-cif}
\end{align}
\end{teiri}
\begin{proof}
\cref{eq:11-fj}：$\{t\le T<t+\Delta,J=j\}\subset\{T\ge t\}$ なので条件付き確率の定義から
$\lambda_j(t)=\lim\frac{\Prob(t\le T<t+\Delta,J=j)}{\Delta\,\Prob(T\ge t)}=f_j(t)/S(t)$（$T$ は連続で $\Prob(T\ge t)=S(t)$）．\\
\cref{eq:11-sum}：イベントは同時に起こらないので，$\{t\le T<t+\Delta\}$ は互いに排反な $\{t\le T<t+\Delta,J=j\}$ $(j=1,\dots,m)$ の和．
両辺を $\Delta\,\Prob(T\ge t)$ で割って極限をとれば $\lambda=\sum\lambda_j$．$S=\exp(-\int_0^t\lambda)$（\cref{thm:09-relations}）に代入する．\\
\cref{eq:11-cif}：$F_j=\int f_j$ に \cref{eq:11-fj} を代入．$\sum_jF_j(t)=\Prob(T\le t)=1-S(t)$．
\end{proof}
CIF は原因 $j$ のハザードだけでなく，$S(u)$ を通じて\textbf{他のすべての原因のハザードに依存する}．
他の原因のハザードが大きいと，原因 $j$ のイベントを起こす前に他のイベントが起きやすいので $F_j$ は小さくなる．

\begin{meidai}[\kakomon{2025}{3}〔2〕]\label{prop:11-Fj}
$m\ge2$ で各原因の確率が正なら，$F_j(\infty)=\Prob(J=j)<1$ であり，$S_j(t)=\Prob(T>t,J=j)$ について $1-F_j(t)>S_j(t)$．
\end{meidai}
\begin{proof}
$\{J=j\}=\{T\le t,J=j\}\cup\{T>t,J=j\}$（排反）より $\Prob(J=j)=F_j(t)+S_j(t)$．
$\Prob(J=j)=1-\sum_{k\ne j}\Prob(J=k)<1$ なので $S_j(t)=\Prob(J=j)-F_j(t)<1-F_j(t)$．
\end{proof}
したがって CIF は「部分分布関数」であり，$t\to\infty$ で1ではなく $\Prob(J=j)$ に収束する．
通常の分布関数のように $1-F_j$ を「生存関数」と解釈することはできない．

\paragraph{例：原因別ハザードが一定の場合}
$\lambda_j(t)=\lambda_j$ なら $S(t)=e^{-(\lambda_1+\lambda_2)t}$，
\[
F_1(t)=\int_0^t\lambda_1e^{-(\lambda_1+\lambda_2)u}\dd u=\frac{\lambda_1}{\lambda_1+\lambda_2}\bigl\{1-e^{-(\lambda_1+\lambda_2)t}\bigr\},\qquad F_1(\infty)=\frac{\lambda_1}{\lambda_1+\lambda_2}.
\]
（独立な $\Ex(\lambda_1)$，$\Ex(\lambda_2)$ の最小値とどちらが先かの同時分布と同じ．）
\paragraph{例：Weibull 型（\kakomon{2025}{3}〔4〕）}
$\lambda_k(t)=\alpha_k\gamma_k(\gamma_kt)^{\alpha_k-1}$ なら $\Lambda_k(t)=(\gamma_kt)^{\alpha_k}$ で $S(t)=\exp\{-(\gamma_1t)^{\alpha_1}-(\gamma_2t)^{\alpha_2}\}$．

\section{推定}\label{sec:11-est}
\subsection{Aalen--Johansen 型推定量}\label{subsec:11-aj}
\Hissu\cref{eq:11-cif}の $\lambda_j(u)\dd u$ を時点ごとの原因別ハザードの推定値 $d_{ji}/n_i$ に，$S(u)$ をその直前の全イベントの KM 推定値に置き換えると
\begin{equation}
\hat F_j(t)=\sum_{i:\,t_i\le t}\frac{d_{ji}}{n_i}\hat S(t_i^-),\label{eq:11-aj}
\end{equation}
ここで $t_i$ は（どれかの原因の）イベント時点，$n_i$ は $t_i$ の直前にどのイベントも起こさず観察中の人数，$d_{ji}$ は $t_i$ での原因 $j$ のイベント数，
$\hat S$ は\textbf{すべての原因をイベントとした} KM 推定量である（\kakomon{2021}{1}〔4〕）．
各時点で「そこまで無事だった確率」×「そこで原因 $j$ が起こる条件付き確率」を足し上げている．
$\sum_j\hat F_j(t)=1-\hat S(t)$ が成り立つ．

\subsection{\texorpdfstring{$1-\mathrm{KM}$}{1−KM} による過大推定}\label{subsec:11-bias}
\Hinshutsu 競合イベントを打ち切りとみなして原因 $j$ だけの KM $\hat S_j^{\mathrm{KM}}$ を計算すると，
各因子は $1-d_{ji}/n_i$ で，$1-\hat S_j^{\mathrm{KM}}(t)$ は $1-\exp\{-\Lambda_j(t)\}$ を推定する．
\begin{meidai}\label{prop:11-overestimate}
$F_j(t)\le1-e^{-\Lambda_j(t)}$．等号は $[0,t]$ 上で（$\lambda_j>0$ となる範囲で）他の原因の累積ハザードが0のときに限る．
\end{meidai}
\begin{proof}
$S(u)=e^{-\Lambda_j(u)}e^{-\sum_{k\ne j}\Lambda_k(u)}\le e^{-\Lambda_j(u)}$ より
$F_j(t)=\int_0^t\lambda_j(u)S(u)\dd u\le\int_0^t\lambda_j(u)e^{-\Lambda_j(u)}\dd u=1-e^{-\Lambda_j(t)}$．
\end{proof}
$1-\mathrm{KM}$ は「競合イベントが起こらない仮想的な世界での原因 $j$ の累積リスク」に相当し，
しかもそれは潜在時間の独立性という検証不能な仮定の下でしか意味を持たない．
競合イベントを起こした人は二度と原因 $j$ のイベントを起こさないのに，打ち切り扱いは
「その人も，リスク集合に残った人と同じ率で将来イベントを起こす」とみなすので，原因 $j$ のイベントを起こさない確率 $1-F_j$ を過小に，累積発生を過大に推定する（\kakomon{2021}{1}〔5〕）．

\begin{figure}[htbp]
\centering
\begin{tikzpicture}
\begin{axis}[width=90mm,height=58mm,xmin=0,xmax=13,ymin=0,ymax=0.65,
  xlabel={時間},ylabel={累積発生},grid=major,grid style={gray!20},tick label style={font=\scriptsize},
  legend style={font=\scriptsize,at={(0.02,0.98)},anchor=north west}]
\addplot[very thick,bkblue,const plot] coordinates {(0,0)(2,0.1)(5,0.2143)(7,0.3286)(9,0.4810)(13,0.4810)};
\addplot[thick,dashed,bkred,const plot] coordinates {(0,0)(2,0.1)(5,0.2286)(7,0.3829)(9,0.5886)(13,0.5886)};
\legend{CIF $\hat F_1(t)$（Aalen--Johansen）,$1-\mathrm{KM}$（原因2を打ち切り扱い）}
\end{axis}
\end{tikzpicture}
\caption{\cref{reidai:11-2}のデータでの累積発生の推定．$1-\mathrm{KM}$ は過大になる．}\label{fig:11-cif}
\end{figure}

\section{回帰モデル}\label{sec:11-reg}
\subsection{原因別ハザードの Cox モデル}
原因 $j$ の原因別ハザードに比例ハザードを仮定する：$\lambda_j(t\mid\bm x)=\lambda_{j0}(t)\exp(\bm\beta_j\transp\bm x)$．
推定は他の原因のイベントを打ち切りとみなした通常の Cox 回帰で行える（原因別ハザードのリスク集合は「どのイベントも起きていない人」だから）．
\textbf{ただし $\bm\beta_j$ は原因別ハザードへの効果であり，CIF への効果ではない}（\kakomon{2025}{3}〔5〕）．
\cref{eq:11-cif}により $F_1(t\mid\bm x)=\int_0^t\lambda_1(u\mid\bm x)\exp\{-\Lambda_1(u\mid\bm x)-\Lambda_2(u\mid\bm x)\}\dd u$ は $\lambda_2$ にも依存し，
$\bm x$ が $\lambda_2$ を大きくするなら，$\bm\beta_1$ の値にかかわらず $F_1$ は小さくなりうる．
原因別ハザードモデルは病因（etiology）の研究，すなわち「まだ生きている人でその原因の発生率が変わるか」に向く．

\subsection{Fine--Gray モデル（部分分布ハザード）}
\Hatten
\begin{teigi}[部分分布ハザード]\label{def:11-subdist}
$\tilde\lambda_j(t)=\lim_{\Delta\to0}\frac{1}{\Delta}\Prob\bigl(t\le T<t+\Delta,J=j\bigm|T\ge t\ \text{または}\ (T<t,\ J\ne j)\bigr)
=\frac{f_j(t)}{1-F_j(t)}=-\frac{\dd}{\dd t}\log\{1-F_j(t)\}$．
\end{teigi}
リスク集合に「すでに他の原因のイベントを起こした人」も含めるのが特徴で，$1-F_j(t)=\exp\{-\int_0^t\tilde\lambda_j\}$ となる．
Fine--Gray モデル \cite{finegray1999} は $\tilde\lambda_1(t\mid\bm x)=\tilde\lambda_{10}(t)\exp(\bm\gamma\transp\bm x)$ とし，
\[
F_1(t\mid\bm x)=1-\{1-F_{10}(t)\}^{\exp(\bm\gamma\transp\bm x)}
\]
から，$\bm\gamma>0$ は「CIF を大きくする」と直接解釈できる．予後予測（ある期間内に原因 $j$ のイベントを経験する確率）に向く．
\begin{supbox}[補足：CIF の群間比較]
2群の CIF の比較には Gray の検定（部分分布ハザードに基づく重み付きログランク型の検定）が用いられる．
原因別ハザードのログランク検定とは帰無仮説が異なる．
\end{supbox}

\section{仮定}\label{sec:11-assump}
\begin{itemize}
\item 打ち切り（追跡不能・試験終了）は独立であること．\textbf{競合イベントは打ち切りではない}．
\item CIF・原因別ハザードは観測データから識別できる．一方，「競合イベントがない世界での原因 $j$ の分布」（潜在時間の周辺分布）は，潜在時間の独立性を仮定しない限り識別できない．
\item 回帰モデル：原因別ハザード（または部分分布ハザード）の比例性．
\end{itemize}

\section{結果の解釈}\label{sec:11-interp}
\begin{itemize}
\item 「5年累積がん死亡率 $\hat F_1(5)=0.20$」：5年以内にがんで亡くなる人の割合の推定値（他因死が起こる現実の世界で）．
\item 高齢者で他因死が多いと，がん死の原因別ハザードは同じでも CIF は低くなる．
\item 治療が原因1の原因別ハザードを下げても，生存が延びた結果として原因2のイベントが増えることがある．両方の原因について結果を示すのが望ましい．
\item estimand の観点では，競合イベントを複合エンドポイント（全死亡）に含める，while alive 戦略で CIF を推定する，などを事前に決める（\cref{sec:02-estimand}）．
\end{itemize}

\section{手法選択}\label{sec:11-choice}
\begin{table}[htbp]
\centering\small
\caption{競合リスクの手法選択}\label{tab:11-choice}
\begin{tabular}{@{}L{55mm}L{80mm}@{}}
\toprule
目的 & 手法\\
\midrule
原因 $j$ の累積発生の記述 & Aalen--Johansen 推定量（$1-\mathrm{KM}$ は不可）\\
CIF の群間比較 & Gray 検定\\
共変量が原因 $j$ の発生率に与える影響（病因） & 原因別ハザードの Cox（競合を打ち切り扱い）\\
共変量が原因 $j$ の累積発生に与える影響（予測） & Fine--Gray モデル\\
全イベントの記述 & 全イベントを1つとした KM・Cox\\
\bottomrule
\end{tabular}
\end{table}

\section{よくある誤り}\label{sec:11-err}
\begin{warnbox}
\begin{itemize}
\item 競合イベントを打ち切りとして $1-\mathrm{KM}$ で累積発生を報告する．
\item Aalen--Johansen 推定量の $\hat S(t^-)$ に，原因 $j$ だけの KM を用いる（全原因の KM を用いる）．
\item 原因別ハザードの HR を「累積発生の比」と解釈する．
\item $F_j(t)\to1$ と考える（極限は $\Prob(J=j)$）．
\end{itemize}
\end{warnbox}

\section{数理統計との接続}\label{sec:11-link}
\begin{linkbox}
\begin{itemize}
\item 独立な $\Ex(\lambda)$ と $\Ex(\mu)$ の最小値 $Z$ と，どちらが最小かの指示変数 $W$ の同時分布（『現代数理統計学の基礎』第4章演習）は，
      一定の原因別ハザードをもつ競合リスクそのものである：$Z\sim\Ex(\lambda+\mu)$，$\Prob(W=1)=\lambda/(\lambda+\mu)$，$Z\indep W$．
\item 最小値のハザードは各ハザードの和（同書第5章演習）が $\lambda=\sum\lambda_j$ に対応する．
\item \kakomonSM{2025}{2} は最小値の分布 $1-\{S(t)\}^n$ を扱った（直列系・最初のイベントの時間）．
\end{itemize}
\end{linkbox}

\section{例題}\label{sec:11-ex}
\begin{reidai}[一定の原因別ハザード]\label{reidai:11-1}
原因1・原因2の原因別ハザードがそれぞれ $0.02$，$0.03$（/年）で一定とする．$e^{-0.5}=0.6065$，$e^{-0.2}=0.8187$ を用いよ．
\begin{enumerate}[label=(\arabic*)]
\item 10年間いずれのイベントも起こさない確率を求めよ．
\item $F_1(10)$，$F_2(10)$，$F_1(\infty)$ を求めよ．
\item 原因2を打ち切りとみなした場合に $1-\mathrm{KM}$ が推定する量 $1-e^{-\Lambda_1(10)}$ を求め，$F_1(10)$ と比べよ．
\end{enumerate}
\end{reidai}

\begin{reidai}[Aalen--Johansen 推定]\label{reidai:11-2}
10例の観察時間とイベントの種類（1：原因1，2：原因2，0：打ち切り）が
$(2,1),(3,2),(4,0),(5,1),(6,2),(7,1),(8,0),(9,1),(10,0),(12,2)$ であった．
\begin{enumerate}[label=(\arabic*)]
\item 各イベント時点の $n_i$，$\hat S(t_i^-)$，$\hat F_1(t)$，$\hat F_2(t)$ を表にせよ．
\item 原因2を打ち切りとみなした $1-\mathrm{KM}$ を $t=9$ で求め，$\hat F_1(9)$ と比べよ．
\item $\hat F_1(12)+\hat F_2(12)+\hat S(12)=1$ を確かめよ．
\end{enumerate}
\end{reidai}

\begin{reidai}[モデルの解釈]\label{reidai:11-3}
高齢がん患者で，治療Aは対照に比べ，がん死の原因別ハザード比が $0.8$，他因死の原因別ハザード比が $1.5$ であった（ともに比例ハザード，原因別ハザードは一定とする）．
対照群のがん死・他因死の原因別ハザードをそれぞれ $0.10$，$0.10$ とする．
\begin{enumerate}[label=(\arabic*)]
\item 各群の $F_{\text{がん}}(\infty)$ を求めよ．
\item 「治療Aはがん死を減らす」という主張を，原因別ハザードと CIF の観点から評価せよ．
\end{enumerate}
\end{reidai}

\section{解答}\label{sec:11-sol}
\begin{kaitou}{reidai:11-1}
(1) $S(10)=e^{-0.05\times10}=e^{-0.5}=0.607$．\\
(2) $F_1(10)=\frac{0.02}{0.05}(1-0.607)=0.4\times0.393=0.157$，$F_2(10)=0.6\times0.393=0.236$，$F_1(\infty)=0.4$．\\
(3) $1-e^{-0.2}=0.181>0.157$．$1-\mathrm{KM}$ は約15\%過大になる．
\end{kaitou}

\begin{kaitou}{reidai:11-2}
(1)
\begin{center}\small
\begin{tabular}{@{}cccccccc@{}}
\toprule
$t_i$ & 種類 & $n_i$ & $\hat S(t_i^-)$ & $\hat S(t_i)$ & $\hat F_1(t_i)$ & $\hat F_2(t_i)$ & $1-\mathrm{KM}_1$\\
\midrule
2 & 1 & 10 & 1.000 & 0.900 & 0.100 & 0.000 & 0.100\\
3 & 2 & 9 & 0.900 & 0.800 & 0.100 & 0.100 & 0.100\\
5 & 1 & 7 & 0.800 & 0.686 & 0.214 & 0.100 & 0.229\\
6 & 2 & 6 & 0.686 & 0.571 & 0.214 & 0.214 & 0.229\\
7 & 1 & 5 & 0.571 & 0.457 & 0.329 & 0.214 & 0.383\\
9 & 1 & 3 & 0.457 & 0.305 & 0.481 & 0.214 & 0.589\\
12 & 2 & 1 & 0.305 & 0.000 & 0.481 & 0.519 & 0.589\\
\bottomrule
\end{tabular}
\end{center}
例：$t=5$ では $\hat F_1=0.1+\frac17\times0.8=0.214$．\\
(2) 原因1だけの KM は $(1-\frac1{10})(1-\frac17)(1-\frac15)(1-\frac13)=0.9\times0.857\times0.8\times0.667=0.411$，
$1-\mathrm{KM}=0.589>\hat F_1(9)=0.481$（\cref{fig:11-cif}）．\\
(3) $0.481+0.519+0=1$．
\end{kaitou}

\begin{kaitou}{reidai:11-3}
(1) 対照：$\frac{0.10}{0.20}=0.50$．治療A：がん $0.08$，他因 $0.15$ で $\frac{0.08}{0.23}=0.348$．\\
(2) がん死の原因別ハザードは20\%低下しており，「生存している人でのがん死の発生率を下げる」という意味では正しい．
CIF も低下しているが，その一部は他因死のハザードが上がって「がん死する前に他因で亡くなる人」が増えた結果である．
他因死の CIF は $0.15/0.23=0.652$（対照0.50）に増え，全死亡のハザードも $0.23>0.20$ と悪化している．
原因別ハザード・CIF・全死亡をあわせて評価すべきで，がん死の CIF の低下だけを根拠に治療を有益とは言えない．
\end{kaitou}

\begin{summarybox}
\begin{itemize}
\item 競合イベントは打ち切りではない．観測は $(T,J)$．
\item $\lambda_j=f_j/S$，$\lambda=\sum\lambda_j$，$S=\exp(-\sum\Lambda_j)$，$F_j(t)=\int_0^t\lambda_jS$，$\sum F_j=1-S$．
\item $F_j(\infty)=\Prob(J=j)<1$，$1-F_j(t)>S_j(t)$．
\item 推定：$\hat F_j(t)=\sum_{t_i\le t}(d_{ji}/n_i)\hat S(t_i^-)$（$\hat S$ は全原因の KM）．
\item $1-\mathrm{KM}$（競合を打ち切り扱い）は $1-e^{-\Lambda_j}\ge F_j$ を推定し過大．
\item 原因別 Cox の $\beta$ はハザードへの効果（病因），Fine--Gray の $\gamma$ は CIF への効果（予測）．
\end{itemize}
\end{summarybox}
